% ============================================================
% SECȚIUNEA 7 — Analiza comparativă
%
% CONVENȚII DE VERIFICARE:
%   [V]  — afirmație verificată din surse publice oficiale sau
%           documentație accesibilă la momentul redactării
%   [NV] — afirmație neverificată; marcată explicit pentru
%           verificare manuală înainte de includere în teză
%
% Platforma proprie (AI Interview Coach) este documentată
% exclusiv din codul sursă; toate afirmațiile despre aceasta
% sunt [V].
% ============================================================

\section{Analiză comparativă}
\label{sec:analiza_comparativa}

\subsection{Platforme analizate}

Comparația acoperă cinci platforme destinate evaluării tehnice a
programatorilor și pregătirii pentru interviuri:

\begin{itemize}
  \item \textbf{AI Interview Coach} — platforma descrisă în această
        lucrare (sursa de referință: codul sursă al proiectului).
  \item \textbf{LeetCode} (\url{leetcode.com}) — platformă de practică
        algoritmică, utilizată preponderent pentru auto-pregătire.
  \item \textbf{HackerRank} (\url{hackerrank.com}) — platformă de
        evaluare tehnică și recrutare, cu suport pentru interviuri live.
  \item \textbf{CoderPad} (\url{coderpad.io}) — mediu de coding
        colaborativ în timp real, focalizat pe interviuri tehnice live.
  \item \textbf{CodeSignal} (\url{codesignal.com}) — platformă de
        evaluare automată a candidaților și practică.
\end{itemize}

\subsection{Tabel comparativ}

\begingroup
\small
\setlength{\tabcolsep}{4pt}
\renewcommand{\arraystretch}{1.35}

\begin{longtable}{p{2.4cm}|p{3.0cm}|p{3.0cm}|p{3.0cm}|p{3.0cm}|p{3.0cm}}
\caption{Comparație între platforme de evaluare tehnică}
\label{tab:comparatie} \\

\hline
\textbf{Dimensiune} &
\textbf{AI Interview Coach} &
\textbf{LeetCode} &
\textbf{HackerRank} &
\textbf{CoderPad} &
\textbf{CodeSignal} \\
\hline
\endfirsthead

\multicolumn{6}{l}{\small\textit{(continuare Tabel~\ref{tab:comparatie})}} \\
\hline
\textbf{Dimensiune} &
\textbf{AI Interview Coach} &
\textbf{LeetCode} &
\textbf{HackerRank} &
\textbf{CoderPad} &
\textbf{CodeSignal} \\
\hline
\endhead

\hline
\multicolumn{6}{r}{\small\textit{continuare pe pagina următoare}} \\
\endfoot

\hline
\endlastfoot

%% ── Feedback AI ──────────────────────────────────────────────────────
\textbf{Feedback AI}
  (format, disponibilitate)
&
  \textbf{[V]}
  Generat de OpenAI (model configurabil);
  4 secțiuni structurate: \emph{Overall}, \emph{Correctness},
  \emph{Code Quality}, \emph{Next Step}.
  Asincron — afișat prin polling după execuție.
  Fallback static local dacă API-ul nu este configurat.
  Inclus în planul de bază (fără nivel premium).
&
  \textbf{[NV]}
  „LeetCode AI" (denumit și AI Mentor / AI Code Review)
  disponibil pentru abonații Premium (din 2023).
  Oferă feedback general pe soluție;
  formatul exact și dacă este asincron sau sincron
  \emph{nu au putut fi verificate}.
&
  \textbf{[NV]}
  Funcționalități AI de code review adăugate în 2023--2024.
  Disponibilitate pe plan și formatul detaliat al
  feedback-ului \emph{nu au putut fi verificate}.
&
  \textbf{[NV]}
  Funcționalități AI de rezumat al interviului și evaluare
  raportate în 2023--2024.
  Dacă există feedback AI structurat \emph{pe cod}
  (nu doar pe sesiune) \emph{nu a putut fi verificat}.
&
  \textbf{[NV]}
  Scoring automat bazat pe AI (``Coding Score'');
  feedback textual detaliat pe submisie
  \emph{nu a putut fi verificat}.
\\
\hline

%% ── Hint-uri ─────────────────────────────────────────────────────────
\textbf{Sistem de hint-uri}
&
  \textbf{[V]}
  Până la 3 niveluri de hint per problemă,
  generate de AI la cerere (endpoint
  \texttt{POST /api/problems/\{id\}/hints}).
  Hint-urile sunt progresive și stocate pe durata sesiunii.
&
  \textbf{[V]}
  1--3 hint-uri statice per problemă, predefinite de autori.
  Soluții editoriale text disponibile.
  Fără generare AI la cerere.
&
  \textbf{[NV]}
  Soluții editoriale și discuții comunitare confirmate.
  Sistem formal de hint-uri progresive
  \emph{nu a putut fi verificat}.
&
  \textbf{[NV]}
  Sistem de hint-uri formal
  \emph{nu a putut fi verificat}.
  Platforma este focalizată pe coding colaborativ live,
  nu pe parcursuri de practică structurată.
&
  \textbf{[NV]}
  Sistem de hint-uri formal
  \emph{nu a putut fi verificat}.
\\
\hline

%% ── Mod interviu vs. practică ────────────────────────────────────────
\textbf{Mod interviu vs.\ practică}
&
  \textbf{[V]}
  \textbf{Ambele} — sesiuni de interviu cu token de acces
  unic, timer, scor agregat și vizualizare de către
  intervievator; mod practică standalone fără sesiune.
  Roluri separate: Candidate, Interviewer, Admin.
&
  \textbf{[V]}
  Predominant practică algoritmică.
  ``Mock Interview'' disponibil.
  ``LeetCode for Companies'' (evaluare candidați) — \textbf{[NV]}
  detalii despre modul interviu gestionat de recrutor
  \emph{nu au putut fi verificate}.
&
  \textbf{[V]}
  \textbf{Ambele} — HackerRank for Work (sesiuni de interviu
  tehnic gestionate de recrutor) și provocări de practică
  publice.
&
  \textbf{[V]}
  Focalizat pe \textbf{interviu live} (coding colaborativ
  în timp real, sesiuni cu partajare de ecran).
  CoderPad Practice (practică asincronă) confirmat,
  dar ca produs secundar.
&
  \textbf{[V]}
  \textbf{Ambele} — Assessments (evaluare automată a
  candidaților) și Arcade / practice (practică).
\\
\hline

%% ── Limbaje suportate ────────────────────────────────────────────────
\textbf{Limbaje de programare}
&
  \textbf{[V]}
  \textbf{3 limbaje}: C\#, Python, C++.
  Executate în sandbox local (DotnetCodeExecutor).
  Extensibil prin configurare, dar nu implicit.
&
  \textbf{[V]}
  \textbf{$\approx$\,20 limbaje}: Python, Java, C++, C\#,
  JavaScript/TypeScript, Go, Rust, Swift, Kotlin etc.
  (numărul exact variază; conform documentației publice).
&
  \textbf{[V]}
  \textbf{$>\,30$ limbaje}, inclusiv limbaje mai puțin
  frecvente (Perl, Scala, Erlang etc.);
  confirmată ca unul dintre punctele forte ale platformei.
&
  \textbf{[V]}
  \textbf{$>\,30$ limbaje}, inclusiv suport pentru
  framework-uri frontend (React, Vue).
  Unul dintre punctele forte diferențiatoare față de
  platformele algoritmice.
&
  \textbf{[NV]}
  Suportă multiple limbaje; numărul exact
  \emph{nu a putut fi verificat}.
\\
\hline

%% ── Self-hosting ─────────────────────────────────────────────────────
\textbf{Self-hosting / open-source}
&
  \textbf{[V]}
  \textbf{Da} — infrastructură proprie;
  .NET 10 + PostgreSQL; fără dependențe de licențe terță
  parte pentru rularea de bază.
  Cheie OpenAI opțională (fallback static fără ea).
  Codul sursă controlat privat; nu este publicat open-source
  în prezent.
&
  \textbf{[V]}
  \textbf{Nu} — exclusiv SaaS.
  Nu există opțiune de deployment on-premises confirmată.
&
  \textbf{[V]}
  \textbf{Nu} — exclusiv SaaS.
  Plan Enterprise disponibil, dar fără on-premises confirmat.
&
  \textbf{[NV]}
  Predominant SaaS.
  Dacă planul Enterprise include opțiune on-premises
  \emph{nu a putut fi verificat}.
&
  \textbf{[V]}
  \textbf{Nu} — exclusiv SaaS.
\\
\hline

\end{longtable}
\endgroup

\subsection{Interpretare și poziționare}

\paragraph{Avantajele distinctive ale platformei propuse}
Față de competitorii analizați, AI Interview Coach se diferențiază prin:

\begin{enumerate}
  \item \textbf{Feedback AI structurat și inclus implicit} — celelalte
        platforme fie nu oferă feedback AI pe cod (CoderPad, CodeSignal),
        fie îl restricționează la abonați premium (LeetCode), fie au
        detalii neverificate (HackerRank). Platforma propusă include
        feedback AI cu 4 secțiuni precise pentru orice utilizator, fără
        nivel premium, cu fallback static dacă serviciul extern nu este
        disponibil.

  \item \textbf{Hint-uri AI progresive, la cerere} — LeetCode oferă
        hint-uri statice predefinite de autor; celelalte platforme nu au
        un sistem de hint-uri confirmat. Platforma propusă generează
        hint-uri la cerere în până la 3 niveluri, adaptate codului
        candidatului.

  \item \textbf{Integrare nativă intervievator-candidat} — HackerRank și
        CoderPad oferă funcționalitate similară, dar ca SaaS. Platforma
        propusă permite unui intervievator să creeze un interviu cu
        probleme proprii, să emită un token de acces și să urmărească
        sesiunile candidaților, totul self-hosted.

  \item \textbf{Self-hosting complet} — niciuna dintre platformele
        concurente analizate nu confirmă o opțiune de deployment
        on-premises. Platforma propusă rulează integral pe infrastructura
        proprie, fără date trimise la terți (cu excepția opțională a API-ului
        OpenAI).
\end{enumerate}

\paragraph{Limitări față de competitori}
\begin{itemize}
  \item \textbf{Limbaje suportate}: 3 față de 20--30+. HackerRank și
        CoderPad acoperă semnificativ mai multe limbaje, inclusiv
        framework-uri frontend.
  \item \textbf{Ecosistem de probleme}: LeetCode și HackerRank dispun
        de cataloage cu mii de probleme; platforma propusă depinde de
        problemele create de intervievatori, plus un set de pornire.
  \item \textbf{Funcționalități de colaborare live}: CoderPad oferă
        coding colaborativ sincron (editor partajat în timp real);
        platforma propusă nu are această funcționalitate.
\end{itemize}

% ─────────────────────────────────────────────────────────────────────────

\section{Schelet de bibliografie}
\label{sec:bibliografie_schelet}

Această secțiune listează \emph{categoriile de surse} necesare pentru
fundamentarea teoretică a lucrării și, unde există consens academic sau
industrial, identifică lucrările canonice de referință.
\textbf{Nu sunt incluse referințe fabricate; fiecare sursă de mai jos
este reală și verificabilă.}
Sursele marcate cu \emph{[de identificat]} necesită căutare activă
pentru a găsi cele mai relevante lucrări actuale.

\subsection*{A. Arhitectura software}

\begin{description}
  \item[Clean Architecture]
    Robert C.\ Martin, \emph{Clean Architecture: A Craftsman's Guide to
    Software Structure and Design}, Prentice Hall, 2017.
    Lucrarea canonică pentru principiile de stratificare folosite în
    proiect (inversiunea dependențelor, separarea responsabilităților).

  \item[Domain-Driven Design]
    Eric Evans, \emph{Domain-Driven Design: Tackling Complexity in the
    Heart of Software}, Addison-Wesley, 2003.
    Sursă pentru conceptele de entitate, aggregate root, repository și
    serviciu de domeniu.

  \item[Patterns of Enterprise Application Architecture]
    Martin Fowler, \emph{Patterns of Enterprise Application Architecture},
    Addison-Wesley, 2002.
    Referință pentru patternuri de infrastructură: Repository, Unit of Work,
    Service Layer (implementate prin EF Core și stratul Application).

  \item[Dependency Injection în .NET]
    Mark Seemann, \emph{Dependency Injection Principles, Practices, and
    Patterns}, Manning, 2019.
    Justificarea alegerilor de lifetime (singleton, scoped, transient) din
    \texttt{ServiceCollectionExtensions}.
\end{description}

\subsection*{B. Modele lingvistice mari și evaluare automată a codului}

\begin{description}
  \item[Mecanismul Transformer (fundament teoretic)]
    Vaswani et al., ``Attention Is All You Need'',
    \emph{Advances in Neural Information Processing Systems}, 2017.
    Arhitectura de bază a modelelor de limbaj utilizate prin API-ul OpenAI.

  \item[Modele LLM pentru generarea și evaluarea codului — \emph{[de identificat]}]
    Căutați lucrări recente (2021--2025) despre: Codex (Chen et al., OpenAI),
    Code Llama (Meta AI), GPT-4 Technical Report (OpenAI, 2023),
    DeepMind AlphaCode.
    Aceste lucrări prezintă benchmark-uri de generare de cod și sunt
    relevante pentru secțiunea de feedback AI.

  \item[AI în educație și sisteme tutoriale inteligente — \emph{[de identificat]}]
    Căutați lucrări în conferințele ACM SIGCSE, ICER sau AIED (Artificial
    Intelligence in Education) despre: automated feedback for programming
    assignments, intelligent tutoring systems for code.

  \item[Evaluare automată a submisiilor de cod (online judges) — \emph{[de identificat]}]
    Căutați lucrări despre arhitecturi de online judge
    (de ex.\ sisteme sandboxate, compararea output-urilor, execuție izolată).
\end{description}

\subsection*{C. Documentații tehnice oficiale (surse primare, stabile)}

\begin{description}
  \item[ASP.NET Core]
    Microsoft, \emph{ASP.NET Core documentation},
    \url{https://learn.microsoft.com/aspnet/core}.
    Referință pentru Minimal Controllers, middleware pipeline, autentificare
    JWT, rate limiting.

  \item[Entity Framework Core]
    Microsoft, \emph{Entity Framework Core documentation},
    \url{https://learn.microsoft.com/ef/core}.
    Referință pentru migrări, relații, InMemory provider (teste).

  \item[Angular]
    Google, \emph{Angular documentation},
    \url{https://angular.dev}.
    Referință pentru Signals, HttpInterceptorFn, CanActivateFn, DestroyRef.

  \item[OpenAI API]
    OpenAI, \emph{OpenAI Platform documentation},
    \url{https://platform.openai.com/docs}.
    Referință pentru endpoint-ul \texttt{/v1/responses}, structura răspunsului,
    câmpul \texttt{max\_output\_tokens}.

  \item[.NET 10]
    Microsoft, \emph{.NET documentation},
    \url{https://learn.microsoft.com/dotnet}.
    Referință pentru \texttt{System.Threading.Channels},
    \texttt{System.Security.Cryptography},
    \texttt{KeyDerivation.Pbkdf2}.
\end{description}

\subsection*{D. Securitate}

\begin{description}
  \item[JWT — Standard]
    M.\ Jones, J.\ Bradley, N.\ Sakimura,
    \emph{RFC 7519: JSON Web Token (JWT)},
    IETF, 2015.
    Specificația formală a formatului tokenurilor folosite în aplicație.

  \item[PBKDF2 — Standard]
    B.\ Kaliski, \emph{RFC 2898: PKCS \#5: Password-Based Cryptography
    Specification Version 2.0}, IETF, 2000.
    Specificația algoritmului de hash al parolelor (PBKDF2-HMAC-SHA256).

  \item[OWASP Top 10]
    OWASP Foundation, \emph{OWASP Top 10},
    \url{https://owasp.org/www-project-top-ten/}.
    Referință pentru vulnerabilitățile adresate: injection, broken access
    control, insecure design, security misconfiguration.
\end{description}

\subsection*{E. Testare software}

\begin{description}
  \item[Unit Testing]
    Roy Osherove, \emph{The Art of Unit Testing} (ed.\ a 3-a),
    Manning, 2023.
    Justificare pentru structura testelor (Arrange-Act-Assert),
    utilizarea InMemory vs.\ mock-uri reale, izolarea serviciilor.

  \item[xUnit.net]
    xUnit.net contributors, \emph{xUnit.net documentation},
    \url{https://xunit.net}.
    Cadrul de testare utilizat în proiect (\texttt{AIInterviewCoach.Tests}).
\end{description}

\subsection*{F. Platforme comparate — surse de verificat}

Afirmațiile marcate \textbf{[NV]} din tabelul comparativ necesită
verificare în:

\begin{itemize}
  \item Paginile oficiale de funcționalități și prețuri ale fiecărei
        platforme (verificate la data citării).
  \item Bloguri de produs și note de versiune oficiale (changelogs)
        pentru funcționalitățile AI adăugate recent (2023--2025).
  \item Documentații tehnice publice (dacă există API-uri publice).
\end{itemize}

\noindent
\textbf{Atenție:} funcționalitățile SaaS se modifică frecvent; orice
comparație trebuie însoțită de data la care a fost efectuată verificarea.
